../cictikz/src/cictikz/data/tex/cictikz_preamble.tex

% Integral and differential non-linearity of a 5-bit DAC.
%
% The DAC is deliberately imperfect, with gain error, curvature and a
% little random mismatch, because a perfect one has nothing to show.
%
% INL, in the middle, is the distance from the best straight line, and it
% measures whether the transfer curve is straight. DNL, at the bottom, is
% the error in each individual step, and it measures whether any step is
% missing. They answer different questions: a converter can have small DNL
% and large INL if it bends smoothly, and small INL with large DNL if one
% step is wrong and the rest compensate.
%
% The mismatch is drawn with a fixed seed, 4, so this figure is the
% same every time it is built. It used to come from an unseeded generator,
% which meant the figure in the book could not be reproduced.
%
% Generated by py/tikzplot.py. Edit the script in ex/ that
% produced it and re-run that, not this file.

\begin{tikzpicture}[thick]

\begin{groupplot}[
  group style={group size=1 by 3, horizontal sep=1.7cm, vertical sep=1.5cm},
  scale only axis,
  grid=both,
  grid style={gray!22, very thin},
  tick align=outside,
  tick label style={font=\small},
  label style={font=\small},
  title style={font=\small, yshift=-1mm},
  legend cell align=left,
  legend style={font=\small, draw=gray!50, fill=white, fill opacity=0.85, text opacity=1}
]
\nextgroupplot[
  width=9.5cm,
  height=3.0cm,
  ylabel={Output voltage [V]},
  legend pos=north west,
  xticklabels={}
]
\addplot[blue, very thick, mark=none] coordinates {(0,-0.015829) (1,0.014657) (2,0.045701) (3,0.076002) (4,0.10605) (5,0.1375) (6,0.16843) (7,0.19998) (8,0.23094) (9,0.26325) (10,0.29519) (11,0.3278) (12,0.35984) (13,0.39259) (14,0.42493) (15,0.45935) (16,0.49159) (17,0.52638) (18,0.56009) (19,0.59441) (20,0.62931) (21,0.66449) (22,0.70036) (23,0.73612) (24,0.77292) (25,0.80859) (26,0.84627) (27,0.88352) (28,0.922) (29,0.95949) (30,0.9989) (31,1.0384)};
\addlegendentry{DAC output}
\addplot[armygreen, very thick, mark=none] coordinates {(0,-0.036237) (1,-0.0024086) (2,0.03142) (3,0.065248) (4,0.099076) (5,0.1329) (6,0.16673) (7,0.20056) (8,0.23439) (9,0.26822) (10,0.30205) (11,0.33587) (12,0.3697) (13,0.40353) (14,0.43736) (15,0.47119) (16,0.50502) (17,0.53884) (18,0.57267) (19,0.6065) (20,0.64033) (21,0.67416) (22,0.70798) (23,0.74181) (24,0.77564) (25,0.80947) (26,0.8433) (27,0.87713) (28,0.91095) (29,0.94478) (30,0.97861) (31,1.0124)};
\addlegendentry{Best fit line}
\nextgroupplot[
  width=9.5cm,
  height=3.0cm,
  ylabel={INL [LSB]},
  xticklabels={}
]
\addplot[blue, very thick, mark=none] coordinates {(0,0.65306) (1,0.54609) (2,0.45701) (3,0.34413) (4,0.22331) (5,0.14721) (6,0.054317) (7,-0.018574) (8,-0.11028) (9,-0.15909) (10,-0.21926) (11,-0.25839) (12,-0.31558) (13,-0.34994) (14,-0.39761) (15,-0.37876) (16,-0.42964) (17,-0.3989) (18,-0.40252) (19,-0.38683) (20,-0.35253) (21,-0.3094) (22,-0.24399) (23,-0.18204) (24,-0.08694) (25,-0.028189) (26,0.095012) (27,0.20448) (28,0.35332) (29,0.4707) (30,0.64928) (31,0.83056)};
\nextgroupplot[
  width=9.5cm,
  height=3.0cm,
  xlabel={Digital code},
  ylabel={DNL [LSB]}
]
\addplot[red, very thick, mark=none] coordinates {(0,-0.024465) (1,-0.006579) (2,-0.03037) (3,-0.038325) (4,0.0064132) (5,-0.010394) (6,0.009613) (7,-0.0092064) (8,0.033697) (9,0.022336) (10,0.043371) (11,0.025313) (12,0.048143) (13,0.034834) (14,0.10136) (15,0.031619) (16,0.11324) (17,0.078884) (18,0.098191) (19,0.11681) (20,0.12564) (21,0.14791) (22,0.14445) (23,0.17761) (24,0.14126) (25,0.2057) (26,0.19198) (27,0.23134) (28,0.19988) (29,0.26109) (30,0.26378)};
\end{groupplot}

\end{tikzpicture}
\end{document}
